\chapter{Discussion}
\label{ch:discussion}

The results of the previous chapter are a collection of specific findings: a step size that does not transfer, an annealing artifact, a Metropolis--Hastings breakdown, a gradient that carries no information, a likelihood trap, three GFlowNet failure modes, and a constraint that does not steer. This chapter draws them together into a single account, states explicitly which research question each part of the results answers and where, connects the findings back to the related work of Chapter~\ref{ch:related}, considers what the account does and does not license us to conclude, and sets out the experiment that would falsify it.

\section{Answers to the Research Questions}
\label{sec:disc-rqs}

It is useful to state plainly which research question each result answers and in which section of Chapter~\ref{ch:results} the answer is established.

\textbf{RQ1}, whether the frozen likelihood can be sampled effectively with faithful Langevin dynamics on the discrete manifold and if not why, is answered in the negative across Sections~\ref{sec:results-stepsize}, \ref{sec:results-fallacy}, and \ref{sec:results-linradius}. Section~\ref{sec:results-stepsize} shows that the sampler cannot even be made to move without model-specific calibration; Section~\ref{sec:results-fallacy} shows, through the direction-versus-magnitude ablation, that following the true gradient is indistinguishable from following a random direction of equal norm; and Section~\ref{sec:results-linradius} supplies the mechanism, namely that the linearization radius of the surrogate is smaller than the smallest admissible discrete step, so the gradient is inapplicable rather than merely imprecise.

\textbf{RQ2}, the role of the Metropolis--Hastings correction in each sampler, is answered in Sections~\ref{sec:results-quench} and \ref{sec:results-mh}. For the discrete sampler the correction is beneficial and removes the quenching artifact (Section~\ref{sec:results-quench}); for the continuous sampler the same correction is catastrophic, and Section~\ref{sec:results-mh} attributes the collapse specifically to the proposal term, connecting it to the non-Lipschitz drift induced by nearest-neighbour projection. The correction does not make the continuous sampler work; it reveals that the continuous sampler was only ever functioning as a biased optimizer.

\textbf{RQ3}, whether amortization with a GFlowNet escapes the failure, is answered in Section~\ref{sec:results-gfn}. It does not: three distinct failure modes appear, and the unifying experiment shows that Langevin sampling on the amortized energy is indistinguishable from sampling on the untuned base, which locates the cause in the energy rather than the search.

\textbf{RQ4}, whether an additive constraint can steer generation on this landscape, is answered in Section~\ref{sec:results-constrained}. It cannot: the steering gains are small and inconsistent in sign, and the failure follows from the same defect established for the fluency energy.

\section{A Single Defect, Reached by Two Routes}
\label{sec:disc-unified}

The organizing claim of this thesis is that all of the failures above are surface manifestations of one underlying defect: the frozen likelihood of an autoregressive language model is not a usable energy function for sampling, because the training objective that produces it never shapes it to be one. Teacher forcing optimizes each local conditional $p_{\text{LM}}(x_t \mid x_{<t})$ to predict the next token given a correct prefix. It places no constraint on how the joint log-likelihood behaves as a function of the sequence when the sequence is varied, and in particular it does not make that function smooth, or well-calibrated as a measure of quality, or navigable by local gradient information. Energy-based decoding borrows this quantity and treats it as though it had those properties. It does not.

The strength of the evidence for this claim rests on the fact that it is reached by two independent routes that share no machinery. The gradient-guided route differentiates the likelihood and finds that the gradient is uninformative (Section~\ref{sec:results-fallacy}), that its surrogate is uncorrelated with the true energy change because the linearization radius is smaller than the smallest discrete step (Section~\ref{sec:results-linradius}), and that the Metropolis--Hastings correction, applied honestly, rejects every useful move because the projected drift is discontinuous (Section~\ref{sec:results-mh}). The amortized route never touches the gradient at all; it treats the likelihood only as a scalar reward, and finds that the reward has no well-behaved optimum where coherent text lies, so the learned policy collapses in whichever direction the reward's implicit incentives push it (Section~\ref{sec:results-gfn}). That two methods with nothing in common beyond the energy they use should fail in ways that trace back to the same property of that energy is the central argument that the energy, and not the sampler, is the problem. The unifying experiment, in which Langevin sampling on the GFlowNet-tuned energy is indistinguishable from Langevin sampling on the base, closes the loop by showing that changing the model through amortization does not change the sampler's fate.

The likelihood trap ties the two routes together at the level of intuition. Because low energy is degenerate text, the goal that both methods pursue is the wrong goal. The gradient-guided samplers cannot reach it, and that is fortunate, because reaching it, as the quenching phase and the GFlowNet reward hacking both show, produces bad text. The methods are, in a sense, protected from their objective by their inability to optimize it, and the moments when they do optimize it effectively are exactly the moments when their output degenerates.

\section{Connections to the Related Work}
\label{sec:disc-related}

The findings of this thesis stand in a definite relationship to the literature reviewed in Chapter~\ref{ch:related}, and it is worth making that relationship explicit rather than leaving it implicit.

The likelihood trap measured in Section~\ref{sec:results-trap} is a direct, quantitative reproduction on our own models of the neural-text-degeneration phenomenon that \citet{holtzman2020curious} identified: that the most probable sequences under a neural language model are degenerate and that maximization-based decoding produces worse text than sampling. What Chapter~\ref{ch:related} presented as a known hazard of decoding, this thesis re-derives as a structural obstacle for energy-based sampling specifically, because it means the target of the sampling procedure is itself undesirable. Our contribution is to connect that degeneration to the sampling objective and to measure the monotone likelihood-repetition relationship directly on the models under study.

The anisotropy results of Section~\ref{sec:results-aniso} confirm, on GPT-2 Large and Llama-3, the representation-degeneration and embedding-anisotropy phenomena documented by \citet{gao2019representation} and \citet{ethayarajh2019contextual}. Chapter~\ref{ch:related} raised anisotropy as a reason to distrust Euclidean distance in embedding space; the measurements here quantify that distrust and, more importantly, show a consequence those works did not consider, namely that the anisotropy scale differs enough between models to make a single step-size schedule non-transferable, which is the concrete cause of the calibration failure in Section~\ref{sec:results-stepsize}.

The Metropolis--Hastings breakdown of Section~\ref{sec:results-mh} is the empirical counterpart to a gap noted in Section~\ref{sec:rel-energy}: that the energy-based decoding literature frequently omits the correction, and where it is included, the difficulty of computing the reverse proposal on a projected landscape is seldom examined. This thesis takes the correction seriously, implements it faithfully following the convergence theory of \citet{roberts1996exponential}, and measures exactly what the omission was hiding. Methods such as COLD \citep{qin2022cold} and MuCoLa \citep{kumar2022gradient}, which our continuous sampler follows, are in this light best understood as biased optimizers rather than samplers, and their reported successes are successes of optimization under early stopping, not of sampling from the stated energy.

Finally, the GFlowNet results of Section~\ref{sec:results-gfn} qualify the picture presented by \citet{hu2024amortizing}, whose formulation and codebase this thesis builds on. Their work showed that amortized sampling from a frozen-model posterior can yield diverse, transferable samples on tasks where the reward is well-behaved. Our narrower and more critical finding is that when the reward is the raw frozen likelihood used as an energy over free-form sequences, the amortization inherits the pathology of that energy, and the learned policy collapses. The two findings are compatible: the difference lies in whether the reward defines a landscape with a coherent-text optimum, and our contribution is to show that the frozen autoregressive likelihood does not.

\section{What the Results Do and Do Not Show}
\label{sec:disc-scope}

It is important to state the limits of these conclusions precisely. The thesis shows that the frozen likelihood of a \emph{standard autoregressive} language model is a poor energy for gradient-guided and amortized sampling on discrete text, and it explains why in terms of the training objective and the geometry of the token embedding space. It does not show that energy-based controllable generation is impossible in general. It shows that a particular and widely used instantiation of it fails, and it identifies the property responsible.

Several caveats bound the empirical claims. The KL-divergence metric depends on the position of the corrupted token, which is why the central results are argued from the absence of a gap between conditions rather than from small measured differences; a null gap is robust to a noisy metric in a way that a narrow win would not be. The Llama-3 results support only qualitative cross-architecture claims, because the tokenizer and embedding scale differ from the GPT-2 family. The linearization result is that the surrogate is uniformly uninformative at all admissible distances rather than informative up to a sharp radius and useless beyond it; the language of a linearization radius is a useful summary, but the measured correlation is low everywhere the sampler operates rather than crossing a clean threshold. And the constrained-generation extension establishes the absence of reliable steering on this energy, not the impossibility of steering on a better one.

The thesis is also candid about the relationship between its title and its content. The original objective was controllable generation via faithful Langevin dynamics, realized as a two-stage plan in which the first stage would establish that energy-guided sampling on a frozen model is viable and the second would add the constraint term and perform sentiment and toxicity control. The first stage did not hold. Rather than proceed to build control on a foundation the study had shown to be unsound, which would have produced results without interpretable meaning, the work turned to diagnosing why the foundation fails. What it delivers is therefore a systematic diagnosis of the failure of a prerequisite, with the mechanisms identified and the constrained-generation extension included as direct evidence, rather than a demonstration of successful control.

\section{Implications and the Falsifying Experiment}
\label{sec:disc-future}

The most valuable consequence of locating the defect in the training objective is that it makes a sharp, testable prediction. If the reason the frozen autoregressive gradient is useless is that teacher forcing never shapes a score function over sequences, then a model whose training objective \emph{does} shape such a function should not suffer the same fate. \textbf{Diffusion language models} are exactly such models \citep{li2022diffusion, lou2024discrete}. A diffusion model is trained, by construction, to denoise corrupted inputs, and denoising is equivalent to learning the gradient of the log-density of the data, the score function, by the identity of \citet{vincent2011connection} that also underpins score-based generative modeling \citep{song2019generative, ho2020denoising}. Running the identical protocol of this thesis on a diffusion language model therefore constitutes a positive control. If Langevin sampling succeeds there, where it fails on the autoregressive model, the training objective is isolated as the causal variable and the diagnosis of this thesis is confirmed. If it fails there too, the diagnosis is wrong and a deeper cause must be sought. Either outcome is informative, and the experiment is inexpensive because score-based discrete diffusion models sharing the GPT-2 tokenizer are available. This is, in the author's view, the single most important next step, and the design of this thesis was shaped to make it a clean test.

A second and complementary direction is to attack the defect at its source rather than route around it. If the problem is that the likelihood is not a score function, one option is to fine-tune the model so that its sequence log-density is aligned directly with a target energy, using a score-matching objective, which is cheaper than the reinforcement-learning-style training of a GFlowNet and directed precisely at the identified deficiency. Only once an energy landscape is known to be navigable does the additive constraint term of the plug-and-play framework become meaningful, and at that point the original goal of controllable generation becomes reachable. The contribution of this thesis to that longer programme is a map of where the ground gives way, and a falsifiable hypothesis about why.
